% CLAUDE REPORT::WORSTCASE-BRIDGE-0001::2026-11-06
\documentclass{article}
\usepackage{amsmath,amssymb,amsthm}
\newtheorem{prop}{Proposition}
\newtheorem{cor}{Corollary}
\title{Worst-case derivation cost $\overline W_R$: what the bridge $\mathrm{SAT}\in\mathrm P\Rightarrow\overline W_R\le\mathrm{poly}$ actually says}
\author{Scheduled run, session 013MBjVMu8bVZotSGcY1had4}
\date{2026-11-06}
\begin{document}\maketitle
\noindent\textbf{Status: P versus NP remains OPEN.} Nothing below proves the lower bound L or a bridge B3 for any concrete $R$. The results are small, PROVED/EXACT, and serve to locate the obstruction.

\section*{Corpus reviewed}
AGENTS.md; root \texttt{main.tex} ownership markers; governance 2026-10-24\_2346 (ARC/FLMV) and 2026-10-28\_0924 (world ledger); LSPV-1/430 (which already fixes $W_R$, $\overline W_R$ and the separate targets B1, B2, B3, and the Tseitin/Urquhart--Ben-Sasson--Wigderson family); LSPV-1/348, 350, 441 (signed orbit, MUS quotient, paid gate-local boundary) by title and header; Claude world-ledger and T-PORT reports. Not re-read byte-for-byte: the ca.\ 400 other LSPV reports. This run does not touch root \texttt{main.tex} (not an LSPV ownership region).

\section*{Definitions (as in LSPV-1/430)}
$R$ is a sound, complete system of locally checkable rules; $\operatorname{cost}_R(\pi)$ is a size measure computable in time $\mathrm{poly}(|\pi|)$ and polynomially equivalent to $|\pi|$ (number of paid DAG nodes). $W_R(F)=\min_\pi \operatorname{cost}_R(\pi)$ and $\overline W_R(N)=\max_{|\langle F\rangle|\le N,\,F\in\mathrm{UNSAT}}W_R(F)$. The minimum is over all derivations, so strategy discovery is free.

\section*{Result 1: B3 is a statement about $R$ alone}
\begin{prop}\label{p1}
For such an $R$ (a Cook--Reckhow system), $\overline W_R(N)\le \mathrm{poly}(N)$ iff $R$ is a polynomially bounded proof system for UNSAT. Hence
\[
\mathrm{B3}(R):\ \mathrm{SAT}\in\mathrm P\Rightarrow\overline W_R\le\mathrm{poly}
\quad\Longleftrightarrow\quad
\mathrm P=\mathrm{NP}\Rightarrow R\text{ is polynomially bounded}.
\]
\end{prop}
\begin{proof}
The first claim is the definition of polynomially bounded. For the second: $\mathrm{SAT}\in\mathrm P$ gives $\mathrm P=\mathrm{NP}$ (and $\mathrm{NP}=\mathrm{coNP}$). Conversely $\mathrm P=\mathrm{NP}\Rightarrow\mathrm{SAT}\in\mathrm P$ since SAT is NP-complete. \qedhere
\end{proof}

\begin{cor}[Dichotomy for proving B3]\label{c1}
Either $\mathrm P\ne\mathrm{NP}$, in which case B3$(R)$ is vacuously true for every $R$ and carries no information, or $\mathrm P=\mathrm{NP}$, in which case B3$(R)$ asserts that $R$ is polynomially bounded. A proof of B3$(R)$ that avoids proving $\mathrm P\ne\mathrm{NP}$ must therefore derive ``$R$ is polynomially bounded'' from the existence of \emph{some unknown} polynomial-time SAT decider $A$, i.e.\ $R$ must polynomially simulate the decider-trace system $P_A=\{(\text{accepting trace of }A\text{ on }\neg F)\}$ for every candidate $A$ (Krajíček--Pudl\'ak, JSL 1989 style). If $A$'s correctness is not provable inside $R$, $P_A$ is a sound system only under the hypothesis, and no uniform $R$ is known that simulates all of them.
\end{cor}
Consequence: B3 is not ``weaker'' than the pointwise or worst-case simulation theorems in the sense of effort. For Cook--Reckhow $R$, B3$(R)$ follows from $R$ simulating every $P_A$, and existence of a $p$-optimal proof system is open (Krajíček--Pudl\'ak 1989; Pudl\'ak, \emph{Incompleteness in the finite domain}, 2017). B3$(R)$ for an $R$ that is merely proved to be strong, such as Frege or Extended Frege, would imply that a hypothetical SAT algorithm needs no correctness proof beyond what $R$ can already verify. Status of that implication: OPEN. Status of the lower bound L for Frege/EF: OPEN (best known Frege lower bounds are linear).

\section*{Result 2: where a fixed-system argument cannot avoid NP versus coNP}
\begin{prop}\label{p2}
Let $R$ be a Cook--Reckhow system. (i) $\mathrm L(R)$ ($\overline W_R$ superpolynomial) alone implies nothing about $\mathrm P$ vs $\mathrm{NP}$ unless $R$ is $p$-optimal. (ii) $\mathrm L(R)\wedge\mathrm B3(R)\Rightarrow\mathrm P\ne\mathrm{NP}$. (iii) If $R$ is $p$-optimal then $\mathrm L(R)\Rightarrow\mathrm{NP}\ne\mathrm{coNP}$, which is strictly stronger than $\mathrm P\ne\mathrm{NP}$ in the sense that it is not known to follow from it.
\end{prop}
\begin{proof}
(ii) is the contrapositive of Proposition~\ref{p1} combined with L. (i) holds because $\mathrm{NP}=\mathrm{coNP}$ is compatible with all non-optimal $R$ being superpolynomial. (iii) Optimality turns L into a lower bound for every proof system, which by Cook--Reckhow is equivalent to $\mathrm{NP}\ne\mathrm{coNP}$. \qedhere
\end{proof}
So the proposed route has two admissible shapes. Shape~A: a non-optimal, sub-Frege $R$ with L and an \emph{independent} B3 simulation (must be proved from $\mathrm P=\mathrm{NP}$ and nothing else; no example known). Shape~B: an optimal $R$ with L, which is NP${}\ne{}$coNP. The repository's $R_{\mathrm{mix}341}$ is in Shape~A and its B3 is unproved.

\section*{Result 3: hostile-instance control}
Odd-charge Tseitin formulas on expanders (Urquhart 1987; Ben-Sasson--Wigderson 2001) are exponentially hard for Resolution but have polynomial-size derivations under Gaussian elimination over $\mathbb F_2$ (Polynomial Calculus over $\mathbb F_2$, $\mathrm{Res}(\oplus)$, and Frege). Hence any $R$ whose rules include XOR-elimination gives $W_R(\mathrm{Tseitin})=\mathrm{poly}$: the family cannot witness L for a mixed system containing linear algebra. This is CONSISTENT with LSPV-1/430, which treats it as a Resolution-only certificate. Candidate hard families for Shape~A must resist (a) resolution, (b) GF(2) and GF($p$) algebra, (c) cutting planes, (d) Extended-Resolution-style definitions; known unconditional candidates surviving all four: none (Frege/EF lower bounds are open).

\section*{Literature check for this run}
Searched current literature on optimal proof systems and Frege progress. Findings: no resolution of the $p$-optimal question or of Frege/EF lower bounds; an arXiv 2025 result gives $\mathrm{AC}^0[p]$-Frege lower bounds for a DNF family (restricted system, Shape~A-compatible only for far weaker $R$); a preprint claiming NP${}\ne{}$coNP (Lin, arXiv:2406.10476, many revisions) has no evidence of acceptance and is treated as unverified. Classification: Krajíček--Pudl\'ak simulation: \emph{directly transferable theorem} (used in Corollary~\ref{c1}); $\mathrm{AC}^0[p]$-Frege lower bounds: \emph{useful analogy needing a transfer theorem} (restricted system does not contain the mixed rules); Lin preprint: \emph{unverified, not used}.

\section*{Duplication / contradiction audit; what changed}
Proposition~\ref{p1} is a restatement; its value is Corollary~\ref{c1}, which says B3 cannot be treated as a cheap bridge. No contradiction with LSPV-1/430, which already lists B3 as unproved. No stale claim found in the portions read. Adopted: the $W_R$, $\overline W_R$ definitions. Retired: nothing.

\section*{Assignments (next steps for the five roles)}
\begin{itemize}
\item Lower-Bound Lead: state B3 for $R_{\mathrm{mix}341}$ as ``$\mathrm P=\mathrm{NP}\Rightarrow R_{\mathrm{mix}341}$ polybounded'' and test whether any weakening (restrict to deciders provably correct in a named bounded-arithmetic theory) is true.
\item Reduction Builder: certify which $P_A$ traces $R_{\mathrm{mix}341}$ simulates with polynomial overhead (e.g.\ DPLL, CDCL, Gaussian-extended).
\item Certificates: find a family resisting the four classes in Result 3, or prove each candidate fails.
\item Resource Auditor: confirm $\operatorname{cost}_R\approx|\pi|$ polynomially for all paid rules, as Proposition~\ref{p1} needs.
\item Red Team: attack Corollary~\ref{c1} (in particular, whether a non-Cook--Reckhow cost such as work with free nondeterminism is legitimately definable).
\end{itemize}
\textbf{P versus NP remains OPEN. B3, L, B1, B2: OPEN.}
\end{document}
% END_OF_VERIFIED_REPORT::CLAUDE::WORSTCASE-BRIDGE-0001
